\section*{अध्याय \ref{ch:g12:early-earth} --- आरंभिक पृथ्वी और उसका वायुमंडल}

\begin{solution}{exo:g12:early-earth:1}
(a) \omterm{def:g12:early-earth:primitive-atmosphere}{बहिर्गैसन}; (b) \omterm{def:g12:early-earth:stromatolite}{स्ट्रोमैटोलाइट}; (c) \omterm{def:g12:early-earth:bif}{पट्टित लौह संरचना}; (d) \omterm{def:g12:early-earth:red-beds}{लाल
संस्तर}; (e) \omterm{def:g12:early-earth:goe}{महा ऑक्सीकरण घटना}।
\end{solution}

\begin{solution}{exo:g12:early-earth:2}
तरल जल दिखाने वाले ज़र्कॉन (4.4 अरब वर्ष); पिलबारा के पहले \omterm{def:g12:early-earth:stromatolite}{स्ट्रोमैटोलाइट}
(3.43); हैमर्सली की \omterm{def:g12:early-earth:bif}{पट्टित लौह संरचनाएँ} (2.5 से
अधिक); \omterm{def:g12:early-earth:goe}{महा ऑक्सीकरण घटना} का अंत (2.06); अधिकांश \omterm{def:g12:early-earth:bif}{पट्टित लौह
संरचनाओं} का अंत (1.8)।
\end{solution}

\begin{solution}{exo:g12:early-earth:3}
$t = 0$ पर, $L = L_\odot/(1 + 0.4) = L_\odot/1.4 \approx 0.714\,L_\odot$:
लगभग \qty{29}{\percent} मंद, आज की तुलना में लगभग एक-तिहाई कम चमकीला।
\end{solution}

\begin{solution}{exo:g12:early-earth:4}
यह कि तब वायु में लगभग कोई ऑक्सीजन नहीं थी। पाइराइट ऑक्सीजन से नष्ट हो जाता है;
उसके कण नदियों द्वारा लुढ़काए जाकर अक्षत रूप में तभी निक्षेपित हो सकते थे जब
सतह पर वायु और जल में उसकी कमी हो।
\end{solution}

\begin{solution}{exo:g12:early-earth:5}
जलवाष्प, जो संघनित होकर महासागर बनी, और कार्बन डाइऑक्साइड, जो
आज से कहीं अधिक मात्रा में वायु में बनी रही और अपने ग्रीनहाउस प्रभाव से पृथ्वी को
गरम रखा।
\end{solution}

\begin{solution}{exo:g12:early-earth:6}
3 अरब वर्ष पहले: $1 - t/t_\odot = 3/4.567 \approx 0.657$, $L \approx
1/1.263 \approx 0.792\,L_\odot$, $T \approx 255 \times 0.792^{1/4} \approx
\qty{241}{K}$। 2.45 अरब वर्ष पहले: $1 - t/t_\odot \approx 0.536$, $L
\approx 1/1.215 \approx 0.823\,L_\odot$, $T \approx 255 \times 0.823^{1/4}
\approx \qty{243}{K}$।
\end{solution}

\begin{solution}{exo:g12:early-earth:7}
\qty{1}{kg} लोहा $1000/55.85 \approx 17.9$ मोल है; ऑक्सीजन का एक
अणु लोहे के चार आयन लेता है, इसलिए $17.9/4 \approx 4.48$ मोल
\ce{O2}, यानी $4.48 \times 32.00 \approx \qty{143}{g}$ ऑक्सीजन।
\end{solution}

\begin{solution}{exo:g12:early-earth:8}
(1) $2.46 - 2.06 = 0.40$ अरब वर्ष, यानी 400 मिलियन वर्ष। (2) 4.54
से 2.46 अरब वर्ष पहले तक: पृथ्वी के इतिहास का $2.08/4.54 \approx \qty{46}{\percent}$।
(3) लगभग $3.43 - 2.46 \approx 1$ अरब वर्ष: सूक्ष्मजीवों की बनाई
ऑक्सीजन वायु में टिक पाने से पहले लंबे समय तक खपती
रही।
\end{solution}

\begin{solution}{exo:g12:early-earth:9}
\omterm{def:g12:early-earth:bif}{पट्टित लौह संरचना} ऑक्सीजन-रहित एक गहरा समुद्र दिखाती है, जो
घुला लोहा ले जा सकता था, और उसके ऊपर एक उथला, धूप वाला समुद्र जिसमें उसे अवक्षेपित करने लायक ऑक्सीजन थी।
भूमि पर निक्षेपित \omterm{def:g12:early-earth:red-beds}{लाल संस्तर} वायु में और सतह के जल में ऑक्सीजन दिखाता है,
जो उसके कणों के लोहे को जंग लगाने के लिए पर्याप्त थी।
\end{solution}

\begin{solution}{exo:g12:early-earth:10}
शार्क बे, पश्चिमी ऑस्ट्रेलिया के हैमेलिन पूल में, जिसका जल सामान्य समुद्री जल से
दुगुना खारा है। सूक्ष्मजीवों की, मुख्यतः सायनोबैक्टीरिया की, एक चटाई अवसाद के कणों को फँसाकर
बाँधती है और कार्बोनेट को बढ़ाती है; ऊपर एक नई चटाई बढ़ती है,
और इसी तरह, परत-दर-परत।
\end{solution}

\begin{solution}{exo:g12:early-earth:11}
(1) $21 \times 10^{-5} = \qty{0.00021}{\percent}$। (2) $21/10 =
\qty{2.1}{\percent}$।
\end{solution}

\begin{solution}{exo:g12:early-earth:12}
\omterm{def:g10:habitable-planet:equilibrium-temperature}{साम्य तापमान} अपरिवर्तित रहता है यदि $L(1 - A')$, $1 - A$ के बराबर हो:
$1 - A' = 0.706/0.741 \approx 0.953$, इसलिए $A' \approx 0.047$। पृथ्वी को
सूर्य-प्रकाश का \qty{5}{\percent} से कम परावर्तित करना पड़ता, लगभग एक
कृष्णिका की तरह, जो बादलों और कुछ बर्फ़ वाला महासागरीय ग्रह नहीं कर सकता।
दूसरा उपाय अधिक प्रबल ग्रीनहाउस प्रभाव है।
\end{solution}

\begin{solution}{exo:g12:early-earth:13}
(1) हिमांक तक पहुँचने के लिए $273 - 237 = \qty{36}{K}$; आज के माध्य तक पहुँचने के लिए
$288 - 237 = \qty{51}{K}$। (2) आज का ग्रीनहाउस प्रभाव $288 - 255 =
\qty{33}{K}$ जोड़ता है: युवा पृथ्वी को अधिक प्रबल प्रभाव चाहिए था। (3) \qty{0.1}{bar}, अकेले
कार्बन डाइऑक्साइड के लिए, आज के पूरे वायुमंडल के दाब का दसवाँ भाग है।
\end{solution}

\begin{solution}{exo:g12:early-earth:14}
(1) 2450 मिलियन वर्ष पहले से पहले चिह्न मौजूद है, इसलिए वायु में आज की
ऑक्सीजन के $10^{-5}$ से कम थी; 2090 मिलियन वर्ष के बाद वह अनुपस्थित है, इसलिए
वायु में अधिक थी। (2) हाँ: \omterm{def:g12:early-earth:goe}{महा ऑक्सीकरण घटना} लगभग 2460 से
2060 मिलियन वर्ष पहले तक चलती है, वही अंतराल।
\end{solution}

\begin{solution}{exo:g12:early-earth:15}
(1) लगभग $4567 - 60 = 4507$ मिलियन वर्ष पहले। (2) टक्कर के अधिक से अधिक $4507 - 4404 =
103$ मिलियन वर्ष बाद: टक्कर के लगभग सौ मिलियन वर्ष के भीतर पृथ्वी के पास
भूपर्पटी और जल थे।
\end{solution}

\begin{solution}{pb:g12:early-earth:1}
\textbf{1.} $4.567 - 4.0 = 0.567$ अरब वर्ष।

\textbf{2.} $1 - 0.567/4.567 \approx 1 - 0.124 = 0.876$।

\textbf{3.} $L = 1/(1 + 0.4 \times 0.876) = 1/1.350 \approx 0.741$, आज की दीप्ति
का।

\textbf{4.} $0.741 \times 1361 \approx \qty{1008}{W/m^2}$।

\textbf{5.} $1 - 0.741 \approx 0.26$: \qty{26}{\percent} मंद।

\textbf{6.} $T_{\mathrm{eq}} = \left(S(1-A)/(4\sigma)\right)^{1/4}$।

\textbf{7.} $A$ और $\sigma$ स्थिर रहने पर, $T_{\mathrm{eq}}$ एक स्थिरांक गुणा
$S^{1/4}$ है, इसलिए $T'/T = (S'/S)^{1/4}$।

\textbf{8.} $0.741^{1/4} = \sqrt{\sqrt{0.741}} \approx \sqrt{0.861} \approx
0.928$।

\textbf{9.} $255 \times 0.928 \approx \qty{237}{K}$।

\textbf{10.} $237 - 273 \approx \qty{-36}{\celsius}$।

\textbf{11.} यह कि 4.4 अरब वर्ष पहले सतह पर तरल जल, एक महासागर,
था।

\textbf{12.} उथले जल में, जैसे आज जीवित \omterm{def:g12:early-earth:stromatolite}{स्ट्रोमैटोलाइट} बढ़ते हैं।

\textbf{13.} हैमर्सली पर्वतमाला की \omterm{def:g12:early-earth:bif}{पट्टित लौह संरचनाएँ}, जो एक समुद्र-तल पर
निक्षेपित हुईं।

\textbf{14.} हाँ: \qty{237}{K}, \qty{273}{K} से बहुत नीचे है; आज के
वायुमंडल के साथ माध्य सतह जमी हुई होती।

\textbf{15.} चट्टानें ऐसे सूर्य के नीचे तरल जल दिखाती हैं जो आज के वायुमंडल के साथ उसे तरल
रखने के लिए बहुत मंद था: \omterm{def:g12:early-earth:faint-young-sun}{मंद युवा सूर्य की समस्या}।

\textbf{16.} $288 - 255 = \qty{33}{K}$।

\textbf{17.} कम से कम $273 - 237 = \qty{36}{K}$, आज से अधिक।

\textbf{18.} $A' = 1 - 0.706/0.741 \approx 0.047$: लगभग काला ग्रह,
जो संभव नहीं।

\textbf{19.} कार्बन डाइऑक्साइड, आज से कहीं अधिक प्रचुर; आरंभिक सूक्ष्मजीवों द्वारा
बनी मीथेन ने सहायता की होगी।

\textbf{20.} चार अरब वर्ष पहले पृथ्वी का \omterm{def:g10:habitable-planet:equilibrium-temperature}{साम्य तापमान}
\textbf{लगभग 237 K} था, हिमांक से कोई \qty{36}{\celsius} नीचे।
\end{solution}
